“最近睡得怎么样？”

我面前坐着一个约莫四十岁的中年男人，面颊清瘦，戴着眼镜。他把病历放在桌上，没有翻开。

“还好，刚开始来的时候还不太适应，总是失眠，最近感觉休息的已经好多了。”

他点了点头，继续说道，“今天我想和你聊点别的，你住院已经一个月了。”

在精神疗养院中每周都要做一次心理检查，对此我已经习惯了。然而医生虽然说是聊些不同的，然而在我看来其实没什么差别，无非是聊聊你的人生经历，据此来推断你的心理状态。对此我已经毫无感受，尽管我可以使我表达的毫无破绽，或许很快就可以出院去了，然而我还是希望可以在疗养院多待一些时日。

疗养院的日子千篇一律：起床，吃饭，院子里遛两圈，吃药，午睡，再起来，再遛两圈，再吃药，睡觉。刚来的时候觉得每一天都长得没有尽头，后来不知道从哪一天开始，时间忽然变快了。等我想回忆昨天做了什么，却发现昨天和前天一模一样，叠在一起分不清了。

疗养院的人大多是受过一些创伤被送到这里的，不过当然也有例外，然而我同他们交流的时候，大多数时候他们很正常，甚至我觉得比外面的人还要正常，这或许是因为他们承受过疼痛，疼痛总是能够让人更快速地联结在一起。

在我隔壁床上的是一个五十多岁的男人，他很少说话，有一种温和的沉默，我能感觉到他其实有很多话想说，但话到了嘴边就又咽了下去，似乎觉得没有必要去说。然而这种沉默就如我见过的大多数面目模糊的中年人，你很难从他那整洁的衣领和克制的动作里，嗅出任何精神溃散的危险气味。

疗养院时不时会组织一些活动，陶艺是其中之一。那天，我正手忙脚乱地对付一团泥，它在我手心里歪来扭去，怎么都扶不正。偏头一看，隔壁床那个男人面前的陶土已经稳稳地升起，变成了一只薄胎的碗。他很少说话，平时只是安静地坐在床边，衣领永远整整齐齐。此刻他的手托着那团旋转的泥，力道均匀而笃定。

我惊讶于他的天赋，尽管我们二人聊天并不多，但毕竟在同一间房间里住着，我上前搭话：

“周先生，你之前做过这个吗？”

”没有，不过...... 我妻子以前喜欢这个。“ 他摇摇头，视线没有离开那团旋转的泥。

他提到妻子的时候，声调有些变化，我意识到这里一定有些变故发生，我看到他手背上绷起的青筋，终究还是把想问的话咽了下去。

说起小林的故事，倒是简单的多，她是住在我隔壁房间的女大学生。某天深夜，当我从睡梦中惊醒，打着手机的手电筒去走廊尽头的洗手间，经过拐角时，我听到一阵很细的哭声，像是极力憋在掌心里发出的动静。光束扫过去，走廊的阴影里蹲着一个女人。光线晃到她的瞬间，哭声停了。她迅速把脸扭向墙壁。我当时没有看清，后来才知道那是小林。

不过白天的小林看起来没有这些困扰，她喜欢烘焙，养了一只狗，疗养院不能做饭，但她会在周末短暂出院时烤好饼干，带回来分装给病友。分发的时候，她有点像小学生一样兴冲冲地和其他人分享她的烘焙诀窍，比如怎么混合面粉和黄油的比例，放在烤箱的时间等诸如此类的问题。

她的人缘颇为不错，经常有朋友前来探望她，有时隔着玻璃，能看到她的朋友坐在床边，眉头紧锁，双手在空中用力地比划着，像是在激烈地控诉什么。小林没有看她们。她只是坐在床沿，目光越过朋友的肩膀，盯着窗外，或者低着头，一点点撕着自己指甲边缘的倒刺。听到她们因为什么事情争吵起来，有时远望她的朋友表现的很生气，而她则经常神不守舍。

我去食堂吃饭的时候总能碰到一个老人，他没有结婚，没有子女，到晚年颇有些孤独。每次遇到，他都会问我的工作，尽管我每次都说过，可下次遇到他还会问，他喜欢絮絮叨叨地说他年轻时的经历，放在以前我会觉得无聊与不耐烦，但我如今感到一种异常的平静。尽管对于我来说这实在不算是一个有趣的经历，但对他来说已经是人生的全部。重复，不断的重复，当人活得越久之后，会发现人可以讲述的并不是很多。

从诊疗室出来的时候，时间还早。今天下午有剧团来演《三姐妹》，契诃夫的戏。我对话剧一向没有兴趣，总觉得演员的情绪太满了。但在这里，任何打破日程的事都值得一试。我决定吃完饭不午睡了，直接去活动室。

说起来针对话剧，我之前一直没有特别的兴趣，话剧作为和戏曲一样，我似乎有点难以接受一种情绪外放的模式，不过，为了消磨时间，我还是决定去看一下。

下午两点，我走进活动室的时候，里面已经坐了一些人。椅子是折叠的，摆了大概七八排，前面空出来一小块地方当作舞台。没有幕布，没有布景，只有几把椅子，一张桌子。窗帘被拉上了一半，下午的光从另一半照进来，落在空出来的那块地板上。

女大学生坐在靠窗的位置，一只手撑着下巴。隔壁床的男人坐在过道旁边，两只手交叉放在膝盖上，他今天换上了一件深色的外套。后面几排稀稀落落坐着几个我见过但不认识的人。我在靠后的位置坐下。食堂那个老人坐在第一排，他已经睡着了。

台上的灯光亮起来的时候，三个女人已经坐在那里了。一个穿着深色的裙子，坐在桌边。一个手里端着杯子，站在窗前。一个很年轻，坐在角落里，低着头像是在想什么。

在接下来的一个小时里，“莫斯科”这个词，伴随着女演员被不合身的旧戏服勒紧的短促呼吸，像某种劣质的麻醉剂，在这间逼仄的活动室里反复回荡。

这块被半掩的窗帘切割出昏暗光块的场地上，没有幕布，也没有布景。大姐抱怨着日复一日站立教书的疲惫，小妹在电报局的差事里枯萎。她们在发黄的防滑地板上走来走去，努力做出一种专属于剧场的悲伤姿态，徒劳地等待着一趟永远不会在这间疗养院里发车的火车。

男爵走近了。他的声音在低音区微微发颤：“可惜的是一样，只有一样，你不爱我。”

昏黄的脚灯从下方打在伊莉娜的脸上，她微微扬起下巴：“这我自己也没有办法呀，我会做你的太太，我会对你忠实、温顺，只是没有爱，这我可有什么办法呢，我一辈子也没有爱过人。啊，我一直那么梦想着爱情，从老早我就日夜梦想着它了。然而我的心就像一架贵重的钢琴把钥匙丢了似的，所以就永远锁着了。我看你的神色很不宁。”

“我整夜没有睡觉......” 男爵喃喃自语。

我听着公爵与伊莉娜的对话，内心感到一阵绞痛，我认得这使人陶醉、令人痛苦的感情，我曾有过这种感受。伊莉娜使得我想起了理佳，尽管当年理佳从未用如此充满诗意与修饰的语言，但此刻，那把丢失的“钢琴钥匙”却实实在在地锁住了我的咽喉。

我可怜起公爵，真想站起来，鞋跟踩在黏糊糊的地板上，冲台上大喊，“你并不是一个人......”

正当我沉思的时候，我突然听到台上传来一阵哭声。伊莉娜哭着撞向旁边的树木，那个她不爱的男爵，在决斗中被人杀死。

然后大姐奥尔加开口了，她的声音没有哭，也没有颤抖，“音乐演奏得多么动听，真想活下去啊。”

奥尔加还在说。她在说那些死去的人，说他们的痛苦正在变成活着的人的快乐。她说再过一些年，幸福与和平会降临大地。

然而，我看着奥尔加那张面无表情的脸，但她的眼眶中却噙着泪珠。奥尔加真的相信幸福与和平会降临大地吗？

舞台中央的伊莉娜开始抽泣，胸口剧烈地起伏着，口中那些关于“工作”和“明天”的字眼，渐渐融化在一片模糊的哽咽中。

“我们为什么活着，我们为什么痛苦，真希望能懂得啊......”

戏演到最后，军乐声从一只小音箱里响起来，音质很差，像蒙上一层旧布。

灯暗了下来，演员鞠躬。

老人被掌声惊醒，也跟着拍了几下，茫然地望向四周，然后重新闭上了眼睛。小林低着头，像是在哭；隔壁床的中年男人没有鼓掌，一动不动地坐在那里。蓦然间他做了一间我意想不到的事情，他坐在小林的旁边，他将手轻轻放在小林的头，抚摸着她的头发，小林没有说话，反而靠的更近了。

不知是谁叹了一声，很粗的嗓子，叹气声落进刚才那个安静的缝隙里，“唉，你说活着有啥子意思。”

叹息声在活动室里渐渐散去，我站起身，没有再看台上那些收拾道具的演员，转身走出了房间。我回到自己的屋子，走到那张掉漆的桌子前，桌上放着半杯凉透的水，旁边是一本医生让我写日记的空白本子。我想写一些东西。

眼前的痛苦已经足够具体，这使我无力再追问更多意义。
